% Pass 4465: general GQ(s,t) line-signing low-moment and quadrangle-incidence theorem.
\subsection{The apartment-parity trace law for a general $\mathrm{GQ}(s,t)$}
\label{sec:pass4465-general-gq-line-signing}

The $W(3,3)$ identities above are not isolated numerical features.  Let
$\mathcal Q$ be any finite generalized quadrangle of order $(s,t)$, let $G$ be
its point-collinearity graph, and assign one sign $\sigma_\ell\in\{\pm1\}$ to
each geometric line $\ell$.  Every line induces a clique $K_{s+1}$ in $G$ and
every graph edge lies on a unique geometric line.  With
\[
 n=(s+1)(st+1),\qquad
 L=(t+1)(st+1),\qquad
 d=s(t+1),
\]
the line-signed adjacency matrix obeys the parameter-exact trace laws
\begin{align}
 \operatorname{tr}(A_\sigma^2)
   &=nd,\\
 \operatorname{tr}(A_\sigma^3)
   &=6\binom{s+1}{3}\sum_{\ell}\sigma_\ell,\\
 \operatorname{tr}(A_\sigma^4)
   &=nd(2d-1)+24L\binom{s+1}{4}+8W_{\mathcal Q}(\sigma),
 \label{eq:pass4465-general-traces}
\end{align}
where $W_{\mathcal Q}$ is the sum of the products of the four side-line signs
over all induced quadrangles.  The fourth-moment classification is exhaustive:
repeated-edge closed walks give $nd(2d-1)$; a simple $4$-cycle with a chord
contains a triangle and is therefore forced into one geometric line; and every
remaining simple $4$-cycle is an induced generalized-quadrangle apartment using
four distinct lines.

The number of induced quadrangles is
\begin{equation}
 \boxed{
 Q=\frac{n s^2 t^2(t+1)}{8}.}
 \label{eq:pass4465-general-Q}
\end{equation}
If $H$ denotes line--quadrangle incidence and $A^\ast$ the line-collinearity
adjacency matrix (equivalently the point graph of the dual $\mathrm{GQ}(t,s)$),
then one line occurs in
\[
 r=\frac{(s+1)s^2t^2}{2}
\]
quadrangles, two intersecting lines occur together in
\[
 \alpha=s^2t,
\]
and two disjoint lines occur together in
\[
 \beta=\binom{s+1}{2}=\frac{s(s+1)}2.
\]
Consequently
\begin{equation}
 \boxed{
 HH^T=(r-\beta)I+(\alpha-\beta)A^\ast+\beta J.}
 \label{eq:pass4465-general-gram}
\end{equation}
For independent random line signs the distinct quadrangle supports are distinct
quartic Walsh characters, so
\[
 \mathbb E W_{\mathcal Q}=0,
 \qquad
 \operatorname{Var}(W_{\mathcal Q})=Q,
 \qquad
 \mathbb E f_{\mathcal Q}=\frac12,
 \qquad
 \operatorname{sd}(f_{\mathcal Q})=\frac1{2\sqrt Q}.
\]

This theorem also clarifies the recent cross-generalized-quadrangle signing
experiments without reviving the refuted ``coarseness law''.  The dual parameter
sets $(s,t)=(3,9)$ and $(9,3)$ both have
\[
 \boxed{Q=102060}
\]
quadrangle observables, but their line-sign registers have respectively $280$
and $112$ degrees of freedom and their signed cliques have sizes $4$ and $10$.
Thus duality preserves the apartment count while not preserving line-signing
granularity.  The theorem predicts exact low moments and null variances; it does
\emph{not} predict the probability that a random line signing is Ramanujan.
The measured $\mathrm{GQ}(3,9)$ counterexample from the recent signing frontier
therefore remains a genuine boundary on any stronger $(s,t)$-only heuristic.

\begin{center}
\fbox{\parbox{0.92\textwidth}{\small
\textbf{Executable boundary.}  The formulas are generated and the $W(3,3)$
specialisation is rebuilt exactly by
\texttt{analysis/w33\_pass4465\_general\_gq\_line\_signing\_trace.py}, with the
frozen certificate
\texttt{data/PART\_W33\_PASS4465\_GENERAL\_GQ\_LINE\_SIGNING\_TRACE.json}.
The result is a finite generalized-quadrangle theorem.  It does not turn the
line signs into a physical gauge field and does not infer spectral optimality
from the parameters alone.}}
\end{center}

% Pass 4466: correction to the recent Pass-4433 four-cycle frustration normalisation.
\subsection{Correction: legacy ``four-cycle frustration'' versus apartment frustration}
\label{sec:pass4466-four-cycle-audit}

The general apartment theorem exposes a normalization issue in the recent
Pass~4433 signing experiment.  Its helper
\texttt{four\_cycles(A)} was documented as ``Every simple 4-cycle ... counted
once'', but the implementation iterates over every unordered diagonal pair and
every pair of common neighbours.  On $W(3,3)$ it therefore emits
\[
 \boxed{3480\text{ records}}
\]
representing only
\[
 \boxed{1740\text{ underlying simple }C_4\text{s}},
\]
each exactly twice.  The $1740$ split as
\[
 \boxed{1620\text{ induced building apartments}+120\text{ line-internal }C_4\text{s}.}
\]
The second family lies inside the forty geometric $K_4$ lines.  Under a line
signing every edge on such a line has the same sign, so every line-internal
$C_4$ has parity $+1$ and can never be frustrated.

Consequently the historical fraction is not the pure apartment fraction.  For
$W(3,3)$ they are related exactly by
\begin{equation}
 \boxed{
 f_{\rm legacy}=\frac{27}{29}f_{\rm apartment}.}
 \label{eq:pass4466-w33-dilution}
\end{equation}
Its independent-random-line-signing mean is therefore $27/58$, not $1/2$.
This does \emph{not} invalidate the standardized $z$-score quoted in
Pass~4433: both the optimized value and its random baseline are multiplied by
the same positive constant, so the $z$-score is invariant.  The correction is
to the label and raw normalization, not to the separate Ramanujan verdict.

For a general $\mathrm{GQ}(s,t)$ the same audit is parameter-exact.  With
$Q$ induced quadrangles and $L=(t+1)(st+1)$ geometric lines, the number of
line-internal simple four-cycles is
\[
 C_{\rm line}=3L\binom{s+1}{4},
\]
and the legacy helper returns $2(Q+C_{\rm line})$ records.  Hence
\begin{equation}
 \boxed{
 f_{\rm legacy}=\frac{Q}{Q+C_{\rm line}}f_{\rm apartment}.}
 \label{eq:pass4466-general-dilution}
\end{equation}
This matters sharply in the recent dual-pair comparison:
\[
 \mathrm{GQ}(3,9):\quad \frac{f_{\rm legacy}}{f_{\rm apartment}}=\frac{243}{245},
 \qquad
 \mathrm{GQ}(9,3):\quad \frac{f_{\rm legacy}}{f_{\rm apartment}}=\frac{81}{137}.
\]
The same $102060$ apartment count can therefore sit underneath very different
raw legacy normalizations.  Cross-generalized-quadrangle comparisons should use
the induced-apartment observable or apply the exact undilution before comparing
fractions.

\begin{center}
\fbox{\parbox{0.92\textwidth}{\small
\textbf{Correction boundary.}  The executable audit is
\texttt{analysis/w33\_pass4466\_legacy\_four\_cycle\_frustration\_audit.py}.
It corrects cycle enumeration semantics and normalization only.  It does not
withdraw Pass~4433's separate statement that a Ramanujan line signing was found
on $W(3,3)$ while the searched $\mathrm{GQ}(9,3)$ line-signing family remained
above its Ramanujan bound.}}
\end{center}
%
